% --- Archon physics formalization source begin ---
% archon:physics
% archon:covers PhyXMiniProblems/problem_phyx_mini_0670.lean
% archon:source-report reports/phyx_mini/problem_phyx_mini_0670.source.json

\chapter{Physics problem phyx\_mini\_0670}
\label{ch:PhyXMiniProblems_problem_phyx_mini_0670}

\paragraph{Problem source.}
\#\# Physical scenario

A particle starts from rest and accelerates as shown in figure.

\#\# Question

Determine the distance traveled in the first \textbackslash{}( 20.0 \textbackslash{}text\{ s\} \textbackslash{}).

\#\# Answer choices

- A: A: 178m

- B: B: 251m

- C: C: 263m

- D: D: 432m

\#\# Figure caption (auxiliary; use the image as primary evidence)

The image is a graph displaying acceleration over time. Here's a detailed description:

- **Axes**: 
  - The x-axis represents time \textbackslash{}( t \textbackslash{}) in seconds (s), marked at intervals of 5 (from 0 to 20).
  - The y-axis represents acceleration \textbackslash{}( a\_x \textbackslash{}) in meters per second squared (m/s\textbackslash{}(\textasciicircum{}2\textbackslash{})), with markings at intervals of 1 (from -3 to 2).

- **Graph Line**: 
  - The graph is a step function with a thick brown line.
  - From \textbackslash{}( t = 0 \textbackslash{}) to \textbackslash{}( t = 10 \textbackslash{}), the acceleration is constant at 2 m/s\textbackslash{}(\textasciicircum{}2\textbackslash{}).
  - At \textbackslash{}( t = 10 \textbackslash{}), there is a sudden drop to an acceleration of 0 m/s\textbackslash{}(\textasciicircum{}2\textbackslash{}), maintaining until \textbackslash{}( t = 15 \textbackslash{}).
  - At \textbackslash{}( t = 15 \textbackslash{}), the acceleration drops to -3 m/s\textbackslash{}(\textasciicircum{}2\textbackslash{}) and remains constant until \textbackslash{}( t = 20 \textbackslash{}).

- **Grid**: 
  - The background includes a light grid facilitating the reading of values.

- **Labels**: 
  - The y-axis is labeled as \textbackslash{}( a\_x \textbackslash{}) (m/s\textbackslash{}(\textasciicircum{}2\textbackslash{})).
  - The x-axis is labeled as \textbackslash{}( t \textbackslash{}) (s).

The piecewise step function indicates three distinct periods of constant acceleration.

\#\# Dataset metadata

Kinematics; Multi-Formula Reasoning, Physical Model Grounding Reasoning

\paragraph{Recorded answer/context.}
C: C: 263m

\paragraph{Figure/image path.}
/root/proposal\_for\_physic/science-mango/phyx\_mini\_run/phyx\_data/test\_image/670.png

\paragraph{Formalization target.}
create a compiling Lean file with sorry bodies at `PhyXMiniProblems/problem\_phyx\_mini\_0670.lean`. The Lean declarations must preserve the physical quantities, dimensions or dimensional roles, figure labels, governing-law hypotheses, and final relation expressed by this problem.
Use Mathlib/Physlib names found through LeanExplore where available. If a physics API is missing, introduce faithful local abstractions rather than scalar placeholder aliases.

\begin{theorem}[Physics formalization target]
\label{thm:physics:phyx_mini_0670:target}
\lean{PhyXMiniProblems.ProblemPhyXMini0670.problem_phyx_mini_0670}
\uses{def:physics:phyx-mini-0670:phyxminiproblems-problemphyxmini0670-distanceinmeters, def:physics:phyx-mini-0670:phyxminiproblems-problemphyxmini0670-timeinseconds, def:physics:phyx-mini-0670:phyxminiproblems-problemphyxmini0670-particlemotionsetup, def:physics:phyx-mini-0670:phyxminiproblems-problemphyxmini0670-matchesparticlescenario, def:physics:phyx-mini-0670:phyxminiproblems-problemphyxmini0670-matchessuppliedaccelerationfigure, def:physics:phyx-mini-0670:phyxminiproblems-problemphyxmini0670-satisfiesonedimensionalkinematics, def:physics:phyx-mini-0670:phyxminiproblems-problemphyxmini0670-recordeddatasetanswer, def:physics:phyx-mini-0670:phyxminiproblems-problemphyxmini0670-isnearestdisplayeddistance}
The assigned autoformalize agent should translate this physics problem into Lean declarations in the covered file, with theorem and lemma proof bodies written as `by sorry`.
\end{theorem}
\begin{proof}
This is an autoformalization task, not a proof task. Produce statements that can later be proved by the physics prover without weakening the physical model.
\end{proof}

% --- Archon declaration topology begin ---
\section{Declaration topology}
\begin{definition}[velocity Dimension]
  \label{def:physics:phyx-mini-0670:phyxminiproblems-problemphyxmini0670-velocitydimension}
  \lean{PhyXMiniProblems.ProblemPhyXMini0670.velocityDimension}
  The physical dimension of velocity, length divided by time
\end{definition}
\begin{proof}
  This is the declared model definition of velocity Dimension.
\end{proof}
\begin{definition}[acceleration Dimension]
  \label{def:physics:phyx-mini-0670:phyxminiproblems-problemphyxmini0670-accelerationdimension}
  \lean{PhyXMiniProblems.ProblemPhyXMini0670.accelerationDimension}
  The physical dimension of acceleration, length divided by time squared
\end{definition}
\begin{proof}
  This is the declared model definition of acceleration Dimension.
\end{proof}
\begin{definition}[Signed Length Quantity]
  \label{def:physics:phyx-mini-0670:phyxminiproblems-problemphyxmini0670-signedlengthquantity}
  \lean{PhyXMiniProblems.ProblemPhyXMini0670.SignedLengthQuantity}
  A signed one-dimensional position or displacement
\end{definition}
\begin{proof}
  This is the declared model definition of Signed Length Quantity.
\end{proof}
\begin{definition}[Distance Quantity]
  \label{def:physics:phyx-mini-0670:phyxminiproblems-problemphyxmini0670-distancequantity}
  \lean{PhyXMiniProblems.ProblemPhyXMini0670.DistanceQuantity}
  A nonnegative physical path length
\end{definition}
\begin{proof}
  This is the declared model definition of Distance Quantity.
\end{proof}
\begin{definition}[Time Quantity]
  \label{def:physics:phyx-mini-0670:phyxminiproblems-problemphyxmini0670-timequantity}
  \lean{PhyXMiniProblems.ProblemPhyXMini0670.TimeQuantity}
  A nonnegative physical duration
\end{definition}
\begin{proof}
  This is the declared model definition of Time Quantity.
\end{proof}
\begin{definition}[Signed Velocity Quantity]
  \label{def:physics:phyx-mini-0670:phyxminiproblems-problemphyxmini0670-signedvelocityquantity}
  \lean{PhyXMiniProblems.ProblemPhyXMini0670.SignedVelocityQuantity}
  \uses{def:physics:phyx-mini-0670:phyxminiproblems-problemphyxmini0670-velocitydimension}
  A signed one-dimensional velocity component
\end{definition}
\begin{proof}
  This is the declared model definition of Signed Velocity Quantity.
\end{proof}
\begin{definition}[Signed Acceleration Quantity]
  \label{def:physics:phyx-mini-0670:phyxminiproblems-problemphyxmini0670-signedaccelerationquantity}
  \lean{PhyXMiniProblems.ProblemPhyXMini0670.SignedAccelerationQuantity}
  \uses{def:physics:phyx-mini-0670:phyxminiproblems-problemphyxmini0670-accelerationdimension}
  A signed one-dimensional acceleration component
\end{definition}
\begin{proof}
  This is the declared model definition of Signed Acceleration Quantity.
\end{proof}
\begin{definition}[signed Length Readout]
  \label{def:physics:phyx-mini-0670:phyxminiproblems-problemphyxmini0670-signedlengthreadout}
  \lean{PhyXMiniProblems.ProblemPhyXMini0670.signedLengthReadout}
  \uses{def:physics:phyx-mini-0670:phyxminiproblems-problemphyxmini0670-signedlengthquantity}
  Read a signed position in a selected length unit
\end{definition}
\begin{proof}
  This is the declared model definition of signed Length Readout.
\end{proof}
\begin{definition}[distance Readout]
  \label{def:physics:phyx-mini-0670:phyxminiproblems-problemphyxmini0670-distancereadout}
  \lean{PhyXMiniProblems.ProblemPhyXMini0670.distanceReadout}
  \uses{def:physics:phyx-mini-0670:phyxminiproblems-problemphyxmini0670-distancequantity}
  Read a nonnegative path length in a selected length unit
\end{definition}
\begin{proof}
  This is the declared model definition of distance Readout.
\end{proof}
\begin{definition}[time Readout]
  \label{def:physics:phyx-mini-0670:phyxminiproblems-problemphyxmini0670-timereadout}
  \lean{PhyXMiniProblems.ProblemPhyXMini0670.timeReadout}
  \uses{def:physics:phyx-mini-0670:phyxminiproblems-problemphyxmini0670-timequantity}
  Read a duration in a selected time unit
\end{definition}
\begin{proof}
  This is the declared model definition of time Readout.
\end{proof}
\begin{definition}[signed Velocity Readout]
  \label{def:physics:phyx-mini-0670:phyxminiproblems-problemphyxmini0670-signedvelocityreadout}
  \lean{PhyXMiniProblems.ProblemPhyXMini0670.signedVelocityReadout}
  \uses{def:physics:phyx-mini-0670:phyxminiproblems-problemphyxmini0670-signedvelocityquantity}
  Read a signed velocity in coherent length and time units
\end{definition}
\begin{proof}
  This is the declared model definition of signed Velocity Readout.
\end{proof}
\begin{definition}[signed Acceleration Readout]
  \label{def:physics:phyx-mini-0670:phyxminiproblems-problemphyxmini0670-signedaccelerationreadout}
  \lean{PhyXMiniProblems.ProblemPhyXMini0670.signedAccelerationReadout}
  \uses{def:physics:phyx-mini-0670:phyxminiproblems-problemphyxmini0670-signedaccelerationquantity}
  Read a signed acceleration in coherent length and time units
\end{definition}
\begin{proof}
  This is the declared model definition of signed Acceleration Readout.
\end{proof}
\begin{definition}[position In Meters]
  \label{def:physics:phyx-mini-0670:phyxminiproblems-problemphyxmini0670-positioninmeters}
  \lean{PhyXMiniProblems.ProblemPhyXMini0670.positionInMeters}
  \uses{def:physics:phyx-mini-0670:phyxminiproblems-problemphyxmini0670-signedlengthquantity, def:physics:phyx-mini-0670:phyxminiproblems-problemphyxmini0670-signedlengthreadout}
  Position readout in metres
\end{definition}
\begin{proof}
  This is the declared model definition of position In Meters.
\end{proof}
\begin{definition}[distance In Meters]
  \label{def:physics:phyx-mini-0670:phyxminiproblems-problemphyxmini0670-distanceinmeters}
  \lean{PhyXMiniProblems.ProblemPhyXMini0670.distanceInMeters}
  \uses{def:physics:phyx-mini-0670:phyxminiproblems-problemphyxmini0670-distancequantity, def:physics:phyx-mini-0670:phyxminiproblems-problemphyxmini0670-distancereadout}
  Path-length readout in metres
\end{definition}
\begin{proof}
  This is the declared model definition of distance In Meters.
\end{proof}
\begin{definition}[time In Seconds]
  \label{def:physics:phyx-mini-0670:phyxminiproblems-problemphyxmini0670-timeinseconds}
  \lean{PhyXMiniProblems.ProblemPhyXMini0670.timeInSeconds}
  \uses{def:physics:phyx-mini-0670:phyxminiproblems-problemphyxmini0670-timequantity, def:physics:phyx-mini-0670:phyxminiproblems-problemphyxmini0670-timereadout}
  Duration readout in seconds
\end{definition}
\begin{proof}
  This is the declared model definition of time In Seconds.
\end{proof}
\begin{definition}[velocity In Meters Per Second]
  \label{def:physics:phyx-mini-0670:phyxminiproblems-problemphyxmini0670-velocityinmeterspersecond}
  \lean{PhyXMiniProblems.ProblemPhyXMini0670.velocityInMetersPerSecond}
  \uses{def:physics:phyx-mini-0670:phyxminiproblems-problemphyxmini0670-signedvelocityquantity, def:physics:phyx-mini-0670:phyxminiproblems-problemphyxmini0670-signedvelocityreadout}
  Signed velocity readout in metres per second
\end{definition}
\begin{proof}
  This is the declared model definition of velocity In Meters Per Second.
\end{proof}
\begin{definition}[acceleration In Meters Per Second Squared]
  \label{def:physics:phyx-mini-0670:phyxminiproblems-problemphyxmini0670-accelerationinmeterspersecondsquared}
  \lean{PhyXMiniProblems.ProblemPhyXMini0670.accelerationInMetersPerSecondSquared}
  \uses{def:physics:phyx-mini-0670:phyxminiproblems-problemphyxmini0670-signedaccelerationquantity, def:physics:phyx-mini-0670:phyxminiproblems-problemphyxmini0670-signedaccelerationreadout}
  Signed acceleration readout in metres per second squared
\end{definition}
\begin{proof}
  This is the declared model definition of acceleration In Meters Per Second Squared.
\end{proof}
\begin{definition}[Spatial Axis]
  \label{def:physics:phyx-mini-0670:phyxminiproblems-problemphyxmini0670-spatialaxis}
  \lean{PhyXMiniProblems.ProblemPhyXMini0670.SpatialAxis}
  The one-dimensional coordinate axis used in the scenario
\end{definition}
\begin{proof}
  This is the declared model definition of Spatial Axis.
\end{proof}
\begin{definition}[Acceleration Graph Axis Quantity]
  \label{def:physics:phyx-mini-0670:phyxminiproblems-problemphyxmini0670-accelerationgraphaxisquantity}
  \lean{PhyXMiniProblems.ProblemPhyXMini0670.AccelerationGraphAxisQuantity}
  Physical quantities labelling the two graph axes
\end{definition}
\begin{proof}
  This is the declared model definition of Acceleration Graph Axis Quantity.
\end{proof}
\begin{definition}[Acceleration Plateau]
  \label{def:physics:phyx-mini-0670:phyxminiproblems-problemphyxmini0670-accelerationplateau}
  \lean{PhyXMiniProblems.ProblemPhyXMini0670.AccelerationPlateau}
  The three horizontal pieces visible in the acceleration graph
\end{definition}
\begin{proof}
  This is the declared model definition of Acceleration Plateau.
\end{proof}
\begin{definition}[Acceleration Time Figure]
  \label{def:physics:phyx-mini-0670:phyxminiproblems-problemphyxmini0670-accelerationtimefigure}
  \lean{PhyXMiniProblems.ProblemPhyXMini0670.AccelerationTimeFigure}
  \uses{def:physics:phyx-mini-0670:phyxminiproblems-problemphyxmini0670-accelerationgraphaxisquantity, def:physics:phyx-mini-0670:phyxminiproblems-problemphyxmini0670-accelerationplateau}
  The acceleration--time figure. Its plotted function is a scalar readout in the named axis units, while the particle's acceleration remains a dimensionful quantity in ParticleMotionSetup
\end{definition}
\begin{proof}
  This is the declared model definition of Acceleration Time Figure.
\end{proof}
\begin{definition}[Particle Motion Setup]
  \label{def:physics:phyx-mini-0670:phyxminiproblems-problemphyxmini0670-particlemotionsetup}
  \lean{PhyXMiniProblems.ProblemPhyXMini0670.ParticleMotionSetup}
  \uses{def:physics:phyx-mini-0670:phyxminiproblems-problemphyxmini0670-signedlengthquantity, def:physics:phyx-mini-0670:phyxminiproblems-problemphyxmini0670-distancequantity, def:physics:phyx-mini-0670:phyxminiproblems-problemphyxmini0670-timequantity, def:physics:phyx-mini-0670:phyxminiproblems-problemphyxmini0670-signedvelocityquantity, def:physics:phyx-mini-0670:phyxminiproblems-problemphyxmini0670-signedaccelerationquantity, def:physics:phyx-mini-0670:phyxminiproblems-problemphyxmini0670-spatialaxis, def:physics:phyx-mini-0670:phyxminiproblems-problemphyxmini0670-accelerationtimefigure}
  The particle's time-dependent observables. Function arguments are elapsed seconds from the graph origin. Distance traveled is an independent physical observable related to velocity only by the governing kinematic law below
\end{definition}
\begin{proof}
  This is the declared model definition of Particle Motion Setup.
\end{proof}
\begin{definition}[acceleration Profile From Figure]
  \label{def:physics:phyx-mini-0670:phyxminiproblems-problemphyxmini0670-accelerationprofilefromfigure}
  \lean{PhyXMiniProblems.ProblemPhyXMini0670.accelerationProfileFromFigure}
  The right-continuous scalar acceleration profile transcribed from the raster. Values at the idealized jump instants do not affect its time integral. The zero value outside the displayed interval makes no physical claim there
\end{definition}
\begin{proof}
  This is the declared model definition of acceleration Profile From Figure.
\end{proof}
\begin{definition}[Matches Particle Scenario]
  \label{def:physics:phyx-mini-0670:phyxminiproblems-problemphyxmini0670-matchesparticlescenario}
  \lean{PhyXMiniProblems.ProblemPhyXMini0670.MatchesParticleScenario}
  \uses{def:physics:phyx-mini-0670:phyxminiproblems-problemphyxmini0670-timeinseconds, def:physics:phyx-mini-0670:phyxminiproblems-problemphyxmini0670-velocityinmeterspersecond, def:physics:phyx-mini-0670:phyxminiproblems-problemphyxmini0670-particlemotionsetup}
  Written scenario data: motion is along x, starts from rest, and lasts 20 s
\end{definition}
\begin{proof}
  This is the declared model definition of Matches Particle Scenario.
\end{proof}
\begin{definition}[Matches Supplied Acceleration Figure]
  \label{def:physics:phyx-mini-0670:phyxminiproblems-problemphyxmini0670-matchessuppliedaccelerationfigure}
  \lean{PhyXMiniProblems.ProblemPhyXMini0670.MatchesSuppliedAccelerationFigure}
  \uses{def:physics:phyx-mini-0670:phyxminiproblems-problemphyxmini0670-accelerationinmeterspersecondsquared, def:physics:phyx-mini-0670:phyxminiproblems-problemphyxmini0670-particlemotionsetup, def:physics:phyx-mini-0670:phyxminiproblems-problemphyxmini0670-accelerationprofilefromfigure}
  Exact readout of the primary raster, including axis labels and units, grid geometry, the two vertical jumps, and all three acceleration plateaus. It contains no velocity, distance, or answer-choice value
\end{definition}
\begin{proof}
  This is the declared model definition of Matches Supplied Acceleration Figure.
\end{proof}
\begin{definition}[Satisfies One Dimensional Kinematics]
  \label{def:physics:phyx-mini-0670:phyxminiproblems-problemphyxmini0670-satisfiesonedimensionalkinematics}
  \lean{PhyXMiniProblems.ProblemPhyXMini0670.SatisfiesOneDimensionalKinematics}
  \uses{def:physics:phyx-mini-0670:phyxminiproblems-problemphyxmini0670-positioninmeters, def:physics:phyx-mini-0670:phyxminiproblems-problemphyxmini0670-distanceinmeters, def:physics:phyx-mini-0670:phyxminiproblems-problemphyxmini0670-velocityinmeterspersecond, def:physics:phyx-mini-0670:phyxminiproblems-problemphyxmini0670-accelerationinmeterspersecondsquared, def:physics:phyx-mini-0670:phyxminiproblems-problemphyxmini0670-particlemotionsetup}
  One-dimensional kinematics in the SI coordinate chart printed on the graph. The acceleration derivative law excludes the two finite jump times, while velocity continuity rules out impulses. Path length is the integral of speed, not signed velocity. None of these generic laws fixes the requested distance
\end{definition}
\begin{proof}
  This is the declared model definition of Satisfies One Dimensional Kinematics.
\end{proof}
\begin{lemma}[velocity profile from acceleration figure]
  \label{lem:physics:phyx-mini-0670:phyxminiproblems-problemphyxmini0670-velocity-profile-from-acceleration-figure}
  \lean{PhyXMiniProblems.ProblemPhyXMini0670.velocity_profile_from_acceleration_figure}
  \uses{def:physics:phyx-mini-0670:phyxminiproblems-problemphyxmini0670-velocityinmeterspersecond, def:physics:phyx-mini-0670:phyxminiproblems-problemphyxmini0670-particlemotionsetup, def:physics:phyx-mini-0670:phyxminiproblems-problemphyxmini0670-matchesparticlescenario, def:physics:phyx-mini-0670:phyxminiproblems-problemphyxmini0670-matchessuppliedaccelerationfigure, def:physics:phyx-mini-0670:phyxminiproblems-problemphyxmini0670-satisfiesonedimensionalkinematics}
  Integrating the graph from the initial rest condition yields the continuous velocity profile 2t, 20, and 65 - 3t on the three displayed intervals
\end{lemma}
\begin{proof}
  The conclusion follows from the governing relations and figure readouts named in the statement of velocity profile from acceleration figure.
\end{proof}
\begin{lemma}[distance traveled in first twenty seconds]
  \label{lem:physics:phyx-mini-0670:phyxminiproblems-problemphyxmini0670-distance-traveled-in-first-twenty-seconds}
  \lean{PhyXMiniProblems.ProblemPhyXMini0670.distance_traveled_in_first_twenty_seconds}
  \uses{def:physics:phyx-mini-0670:phyxminiproblems-problemphyxmini0670-distanceinmeters, def:physics:phyx-mini-0670:phyxminiproblems-problemphyxmini0670-timeinseconds, def:physics:phyx-mini-0670:phyxminiproblems-problemphyxmini0670-particlemotionsetup, def:physics:phyx-mini-0670:phyxminiproblems-problemphyxmini0670-matchesparticlescenario, def:physics:phyx-mini-0670:phyxminiproblems-problemphyxmini0670-matchessuppliedaccelerationfigure, def:physics:phyx-mini-0670:phyxminiproblems-problemphyxmini0670-satisfiesonedimensionalkinematics}
  The path length accumulated in the first 20 s is exactly 100 + 100 + 125/2 = 525/2 m. The velocity stays nonnegative, ending at 5 m/s, so the speed integral has no reversal correction
\end{lemma}
\begin{proof}
  The conclusion follows from the governing relations and figure readouts named in the statement of distance traveled in first twenty seconds.
\end{proof}
\begin{definition}[Answer Choice]
  \label{def:physics:phyx-mini-0670:phyxminiproblems-problemphyxmini0670-answerchoice}
  \lean{PhyXMiniProblems.ProblemPhyXMini0670.AnswerChoice}
  The four answer labels displayed with the problem
\end{definition}
\begin{proof}
  This is the declared model definition of Answer Choice.
\end{proof}
\begin{definition}[displayed Distance In Meters]
  \label{def:physics:phyx-mini-0670:phyxminiproblems-problemphyxmini0670-displayeddistanceinmeters}
  \lean{PhyXMiniProblems.ProblemPhyXMini0670.displayedDistanceInMeters}
  \uses{def:physics:phyx-mini-0670:phyxminiproblems-problemphyxmini0670-answerchoice}
  Distance readouts in metres printed beside the four answer labels
\end{definition}
\begin{proof}
  This is the declared model definition of displayed Distance In Meters.
\end{proof}
\begin{definition}[recorded Dataset Answer]
  \label{def:physics:phyx-mini-0670:phyxminiproblems-problemphyxmini0670-recordeddatasetanswer}
  \lean{PhyXMiniProblems.ProblemPhyXMini0670.recordedDatasetAnswer}
  \uses{def:physics:phyx-mini-0670:phyxminiproblems-problemphyxmini0670-answerchoice}
  The answer label recorded in the dataset metadata; this is not a premise
\end{definition}
\begin{proof}
  This is the declared model definition of recorded Dataset Answer.
\end{proof}
\begin{definition}[Is Nearest Displayed Distance]
  \label{def:physics:phyx-mini-0670:phyxminiproblems-problemphyxmini0670-isnearestdisplayeddistance}
  \lean{PhyXMiniProblems.ProblemPhyXMini0670.IsNearestDisplayedDistance}
  \uses{def:physics:phyx-mini-0670:phyxminiproblems-problemphyxmini0670-answerchoice, def:physics:phyx-mini-0670:phyxminiproblems-problemphyxmini0670-displayeddistanceinmeters}
  A displayed answer is nearest to an exact distance readout
\end{definition}
\begin{proof}
  This is the declared model definition of Is Nearest Displayed Distance.
\end{proof}
% --- Archon declaration topology end ---
% --- Archon physics formalization source end ---
